\subsection{Culture talk as sociability}
We can begin to characterize the empirical phenomenon will be the focus of analysis via the rather common observation that conversational rituals usually require some sort of \emph{content} for their successful initiation and maintenance \citep{dimaggio87, collins75}. Whenever the primary conversational subject matter animating a given interaction episode moves towards the arts and the symbolic products of the culture industry, we may consider this an instance of  \emph{culture talk}. Following this lead, I define culture talk as \emph{the (more or less) routine deployment of cultural knowledge associated with aesthetic goods as a resource to generate conversation}. In a nutshell: the discussion of ``cultural matters'' with friends and acquaintances.

From a Simmelian perspective, we may consider culture talk one of the primary examples of ``sociability'' in late modern societies. Sociability is social interaction that is not driven by explicit instrumental purposes, but that is keyed to the continuing repetition of episodes of social interaction for their own sake \citep{simmel49}. Sociability thus represents the limiting case---and because seemingly devoid of explicit purpose, the ``play form''---of the more general process of association (coming together from some purpose).  In this respect, sociability is a \emph{form} of social intercourse explicitly dissociated from any sort of goal-directed pursuit. As a social form, sociability is thus partially independent of the particular content that drives it in a given context, making culture talk a fairly general (and generic) phenomenon.  However, this partial independence of form from content does not mean that just any content will do.

\subsubsection{The importance of aesthetic discourse as content}
Recent theorizing on the culture-network link has moved beyond the early formalism of traditional network analysis to propose that specifying the form of association is not enough; instead, particular forms of association are constituted by specific cultural contents that provide meaning to those forms \citep{fuhse09}.  Here I propose that in the contemporary context talk related to (1)  \emph{aesthetic consumption experiences} \citep{holbrook_hirschman82, deighton92, gainer95}, and (2) general knowledge of the accepted discursive conventions governing expressions of aesthetic categorization, evaluation and enjoyment \citep{baumann07, frith98} become one of the most important sources of meaningful content making possible sustained episodes of sociability.

As Simmel noted, sociability cannot be driven by traditional ``practical'' contents (e.g. work, politics, etc.), but requires instead ludic content that serves to simply sustain the focus of attention in the interaction itself \citep{collins75}. The reason for this is that sociability is a kind of social intercourse ``which does not have a strictly economic or business purpose'' but access to which (and restrictions on) is a key marker of the status group boundaries \citep[114]{weber94}---although in other contexts similar kinds of seemingly purposeless content are useful in \emph{relaxing} those same status boundaries \citep{ikegami05}. Building on Simmel, I argue that culture talk is the primary manifestation of sociability in contemporary class-differentiated societies, precisely because conversation related to symbolic goods is the ideal ludic content that fits this form of interaction. In comparison to aesthetic talk other types of conversational content---such as for instance religion or politics---are miserable failures, because they tend to suppress conversation and tend to highlight latent social divisions \citep{eliasoph98}.

Where does aesthetic content come from? Cultural theorists have persuasively argued that media-culture and arts dissemination institutions play a key role in producing an overabundant supply of (sometimes ephemeral) content useful for generating discussion that serves as a resource for conversation in the context of small groups  \citep[453]{fine77}.  The recurrently reproduced micro-orders based on face-to-face (and increasingly electronically mediated) interaction require a constant stream of novel content in order to be recurrently recreated and thus reproduce themselves over time. 

In this respect, the formal and informal system of production of aesthetic goods provide as ``input'' the very content that serves as the scaffolding around which everyday episodes of ``talk'' and sociability are built.  Popular culture scholars have proposed that ``[p]art of the pleasure of popular culture is talking about it; part of its meaning is{\ldots}talk''  \citep[4]{frith98}. Accordingly, civic membership and even ``belonging'' to neighborhoods and communities is partially constituted via talk, and ``[m]uch of this talk is about the mass media and its cultural commodities'' \citep[73]{fiske87}. For instance, Elizabeth Long \citeyearpar{long03} in her study of women's reading groups in Houston, Texas, finds that cultural knowledge and taste become the primary content of social interaction. Consistent with the claim that connects the consumption of aesthetic products and sociability, she finds that reading group members ``tend to press books into service for the meanings that they transmit and the conversations they generate \citeyearpar[73]{long03}. Simlarly, \citet[316-317]{benzecry_collins14} point to the crucial role of culture talk in constituting the aesthetic experience for fanatics of Opera in Buenos Aires. Among Opera fans, talk is almost constant, and this talk revolves around the music, past performances, and tales of their own ``heroic commitment'' to the art form. Through this talk Opera fanatics ritually construct and maintain their relationship to other fans, while at the same time marking the boundary separating true connoisseurs from less serious followers. 

Why is aesthetic talk effective at relationship building and maintenance? As noted in both classic \citep{fine77, dimaggio87} and more recent work \citep{cardon_granjon05, lewis_etal12}, talk about culture and the arts infuses local conversation-based interaction rituals with a common focus of attention, serving as fodder for either deep, engrossing conversations \citep{frith98} or ``for least-common denominator talk'' \citep[443]{dimaggio87}, depending on the pragmatic context. In contrast to material goods which ``are physically present and visible'' and which may serve as barriers for interaction via the creation of invidious distinction, the consumption of aesthetic goods ``\ldots is invisible once it has occurred. This evanescent quality makes artistic experience, described and exploited in conversation, a portable and thus potent medium of interactional exchange.'' \citep[443-444]{dimaggio87}. Because commercial (``popular'') culture and the arts (along with ``fashion, cuisine and sport'') are the closest analog to a common cultural currency in contemporary societies \citep{dimaggio00a}, they have the capacity to be a smooth facilitator of interaction across situational settings, social occasions and social positions\citep{erickson96}. 

Both early \citep{dimaggio87, fine77} and more recent \citep{long03, benzecry_collins14} considerations of the role of culture talk in the creation and maintenance of social relations via consumption and interaction rituals, converges with work in marketing and consumer culture studies \citep{deighton92, gainer95}. According to \citet{gainer95}, art consumers use conversation about the arts to make what could potentially be a purely private experience part of their public interactions with others. Via culture talk, persons con-construct both their own personal identities as art consumers \emph{and} their interpersonal relations with other art aficionados. In this way, ``social relationships are often created and maintained through rituals of reproduction based on consumption which individuals join in shared performances with symbolic meaning.''\citep[254]{gainer95}. Consumption of the arts, as mediated via culture talk, allow persons to pursue interpersonal goals connected to companionship and friendship, build bridges to socially distant acquaintances, generate interaction rituals for work conversation \citep{lizardo13}, establish criteria for co-membership in exclusive social circles \citep{kadushin66}, while at the same time managing and reinforcing individual and community identities \citep[255-257]{gainer95}. 

